\section{What is proved, and what ``in iPRA'' means}
\label{sec:foundations}

The main result is a uniform arithmetic form of the Cayley--Hamilton theorem.
It is accompanied by a finite constructive development of its principal
algebraic consequences. Throughout, a ring is associative, commutative, and
unital unless explicitly stated otherwise. We allow the zero ring. Thus a
polynomial displayed with leading coefficient $1$ is monic in the formal
coefficient sense; its actual degree is $n$ only when $1\ne0$.

\begin{theorem}[Uniform arithmetic Cayley--Hamilton]
\label{thm:arithmetic}
There is a primitive recursive function $\mathrm{CHCheck}:\mathbb N\to
\{0,1\}$, defined explicitly below by integer polynomial arithmetic, for which
intuitionistic primitive recursive arithmetic proves
\[
                 \forall n\;\mathrm{CHCheck}(n)=1.
\]
The assertion says that every coefficient in every entry of
$\det(tI_n-X)\big|_{t=X}$ is zero, where $X=(x_{ij})$ is the generic matrix
over $\mathbb Z[x_{00},\ldots,x_{n-1,n-1}]$. Consequently, specializing the
variables gives $\chi_A(A)=0$ over every commutative ring. For each fixed
primitive recursive presentation of a ring whose ring laws are provable in
iPRA, this specialization is itself an iPRA derivation.
\end{theorem}

The human proof of this theorem is given in Sections~\ref{sec:determinants}
and~\ref{sec:main}, with its induction accounting below. The accompanying
Lean files are a specification and a route to a completed formalization;
they are not a completed Lean kernel verification of this arithmetic
theorem. The executable checks are independent finite corroboration, not
a substitute for the uniform proof.

There are three distinct levels of assertion in this report.
\begin{enumerate}
\item \emph{Arithmetic identities}: equations between specified primitive
recursive functions on natural-number codes. These are the literal iPRA
statements.
\item \emph{Algebraic interpretations}: theorems about arbitrary rings,
modules, or supplied field data obtained from the finite identities. Such
objects are mathematical interpretations, not new sorts silently added to
first-order arithmetic.
\item \emph{Effective analytic versions}: statements about rational
approximations with explicit primitive recursive moduli. These are reduced
to finite inequalities at any requested precision, not quantified
existence claims about arbitrary real numbers or arbitrary analytic
functions.
\end{enumerate}
This distinction is necessary: a short theorem in ordinary Lean with
\texttt{[CommRing R]} does not alone establish an object-language iPRA proof.

\subsection{The base theory and its permitted induction}

We use intuitionistic first-order logic with equality, symbols for the
primitive recursive functions, their defining equations, and induction
for quantifier-free formulas with parameters. A primitive recursive
definition is built from zero, successor, projections, composition, and
primitive recursion. Abbreviations for newly defined primitive recursive
functions are conservative definitional extensions. This is the
presentation underlying the repository interface inspected in
Section~\ref{sec:lean}.

The quantifier-free induction rule has the form
\[
 \frac{\varphi(0,\mathbf a)\qquad
       \varphi(k,\mathbf a)\longrightarrow\varphi(k+1,\mathbf a)}
      {\varphi(k,\mathbf a)},
 \qquad \varphi\text{ quantifier-free}.
\]
Universal closure is taken after the derivation, not inserted into its
induction formula. Equality of natural numbers and every primitive
recursive Boolean test are decidable by their defining recursions. Thus
finite case distinctions and bounded search are legitimate intuitionistically.
No step requires excluded middle for arbitrary formulas, Markov's
principle, choice, an infinite basis construction, or a completeness axiom.

If $P(i,\mathbf a)$ is a decidable primitive recursive predicate, the bounded
assertion $\forall i<b\,P(i,\mathbf a)$ can be replaced by the equation
\[
 \operatorname{all}(b,P,\mathbf a)=1,\qquad
 \operatorname{all}(0,P,\mathbf a)=1,\quad
 \operatorname{all}(k+1,P,\mathbf a)
 =\operatorname{all}(k,P,\mathbf a)\,\mathbf1_{P(k,\mathbf a)}.
\]
Here $P$ is a fixed definition, not a quantified function argument.
An elementary quantifier-free induction proves that the product is $1$
exactly when each test is $1$. Nested bounded assertions are flattened in
the same way. This device will express statements about all entries,
coefficients, permutations, and intermediate stages of a finite calculation.

\subsection{Explicit coding conventions}

One possible list coding uses
\[
 \operatorname{nil}=0,\qquad
 \operatorname{cons}(a,b)=2^a(2b+1).
\]
For a positive code, find its exponent of $2$ by searching up to the code;
the remaining odd factor determines the tail. The tail is smaller than
the code. Repeated decoding for at most the original code many steps
therefore suffices to define length, indexing with default $0$, append,
and bounded updates by primitive recursion. Efficient coding is
irrelevant to proof-theoretic strength; this deliberately simple coding
also shows that every advertised bound is a natural-number bound.

Integers are encoded by $2a$ for $a\ge0$ and $2a+1$ for $-(a+1)$.
Signed addition, multiplication, negation, and comparison are primitive
recursive, and their ring laws have the familiar quantifier-free proofs
from the laws of natural-number arithmetic. A rational is a signed
integer numerator and a positive denominator, optionally reduced by the
Euclidean algorithm. That algorithm can be run for a bound supplied by
its initial arguments, with a stable state after termination. Comparison
and equality use cross multiplication. Invalid input codes receive a
fixed default so that all underlying functions are total.

An $n\times m$ matrix is a list of $nm$ scalar codes, indexed in row-major
order. A univariate polynomial is a coefficient list, interpreted with
zero beyond its length. Matrix multiplication, polynomial convolution,
trace, and powers are bounded sums and products or primitive recursions.
Polynomial coefficients need not be trimmed: degree \emph{bounds}, rather
than a choice of the last nonzero coefficient, suffice throughout the core.

For a completely literal generic identity, use multivariate polynomials
over the coded integers. With $r$ variables and a coordinate exponent
bound $d$, enumerate exponent vectors $\alpha\in\{0,\ldots,d\}^{r}$
by the base-$(d+1)$ number
\[
                    [\alpha]=\sum_{j<r}\alpha_j(d+1)^j.
\]
Store the coefficient array of length $(d+1)^r$. Addition pads to a common
bound. Multiplication increases the bound to the sum of the input bounds
and uses the finite convolution
\[
 (fg)_\alpha=\sum_{\beta,\gamma:\,\beta+\gamma=\alpha}
                       f_\beta g_\gamma.
\]
The sum is implemented over all bounded vector pairs with a Boolean
filter. Thus no unproved normal-form decision procedure is assumed.
Commutativity, associativity, and distributivity follow by finite
reindexing of these convolution sums. At a fixed degree bound, polynomial
equality is equality of two finite arrays.

For dimension $n$, set $r=n^2$ and form the generic matrix $X$. Compute
$c_0,\ldots,c_n$ from the permutation expansion of $\det(tI-X)$, then
compute $\sum_{j=0}^{n}c_jX^j$. The coefficient $c_j$ is supported only on
monomials of total degree $n-j$, and each entry of $X^j$ is supported
only in total degree $j$. Product expansion and induction on $j$ prove
these support assertions, including zero entries. Thus every residual
entry is supported in total degree $n$. No coefficient is discarded by
the coordinate exponent cap $n$, and comparing
the $n^2(n+1)^{n^2}$ integer coefficients suffices. For $n=0$ this is an
empty conjunction, equal to $1$. This algorithm is the promised
$\mathrm{CHCheck}$.

\subsection{Why primitive recursiveness alone is insufficient}
\label{sec:induction}

An algorithm being primitive recursive does not imply that every true
claim about it is provable in iPRA. The proof obligations here are
specific finite loop invariants, not an appeal to that invalid principle.
The following bookkeeping explains how the ordinary notation in the
proofs abbreviates quantifier-free arithmetic reasoning.

Fix an ambient dimension $n$ and the inputs to a claimed identity. All
index sets have explicit bounds: matrix indices below $n$, candidate
permutations among the $n^n$ maps $\{0,\ldots,n-1\}\to\{0,\ldots,n-1\}$,
subsets among $2^n$ bit masks, and polynomial coefficient arrays bounded
by the displayed exponent caps. Use $0^0=1$ for the unique empty map.
The following invariants suffice.
\begin{enumerate}
\item A finite sum after $k$ steps equals the sum of its first $k$ summands;
reindexing follows by inserting each value of a bounded permutation
exactly once. Products have the corresponding invariant.
\item In a sign-reversing involution, the sum over the first $k$ pairs is
zero. Pairs are indexed by the smaller member of their two-element orbit,
selected by a decidable numerical comparison. This proves cancellation
even in characteristic $2$, without cancelling a factor $2$.
\item In convolution, the partial sum accounts for exactly the first $k$
index pairs or triples. Associativity is a reindexing of the same bounded
triple sum, so its packed coefficient invariant is quantifier-free.
\item The adjugate entry identity is the cofactor expansion with a row or
column replaced. Its bookkeeping is the bounded permutation deletion
map described in Section~\ref{sec:determinants}.
\item In the main proof, after processing coefficient $k<n$, the invariant
is exactly
\[
                \sum_{j=0}^{k} c_jX^j=-X^{k+1}B_k.
\]
It is a finite array equality. No induction on the formula ``all matrices
of size $k$ satisfy Cayley--Hamilton'' is used.
\end{enumerate}
These are equations between fixed primitive recursive functions, or
bounded conjunctions of such equations. If a lower-dimensional identity
is needed in a dimension induction, keep $n$ fixed and pack all its
restrictions of sizes at most the current stage into one finite array
invariant. Alternatively, the determinant identities below are derived
directly at the given dimension by finite reindexing.

Integer sizes and list codes do not create an implicit unbounded search.
For the generic calculations, expand a computation without collecting
terms: addition adds term counts and multiplication multiplies them;
the sum of absolute integer coefficients obeys the same bounds.
Primitive recursion through the given computation produces an explicit
bound on every coefficient and code. For other fixed primitive
recursive scalar operations, monotone majorants are constructed along
their definitions. The finite computation is then indexed within those
bounds. This is the standard direct compilation of each stated
invariant into quantifier-free induction, not a uniform total interpreter
for all primitive recursive programs.

Universal polynomial equality can subsequently be evaluated in any
commutative ring: evaluate each finite monomial and sum, using the ring
laws to preserve addition and multiplication. This external
specialization uses no decision of equality in that ring. To make it an
internal iPRA theorem, fix a presentation with primitive recursive
operations and iPRA proofs of its ring laws. Such a presentation uses
either canonical scalar codes with numerical equality, or a primitive
recursive equivalence relation together with proofs of equivalence,
congruence, and the ring laws modulo that relation. Unreduced rational
pairs use the latter option, with cross-multiplication equality.
One must not assert that an
arbitrary program code describes such a ring, or that iPRA decides this.

